% ============================================================
%  J. P. Jacobsen og Vilhelm Møller, "Darwin: hans Liv og hans Lære"
%  Gyldendal, København, 1893
%  Excerpt: "Darwins Trostab" (Darwin's Loss of Faith)
%  Chapter: "Darwins Liv," pp. 38–45
%  TRANSLATION (English)
%
%  Translator's note: "Darwin: hans Liv og hans Lære" (Darwin: His
%  Life and His Doctrine) was published posthumously in 1893 by
%  Gyldendal, collecting J. P. Jacobsen's journal articles on Darwin's
%  major works together with supplementary material added by his friend
%  and editor Vilhelm Møller. The chapter "Darwins Liv" (Darwin's Life)
%  from which this excerpt is taken was written by Møller; Jacobsen's
%  own contributions are the expositions of the Origin of Species and
%  (the first half of) the Descent of Man, which he had translated into
%  Danish in 1871-74. In the present excerpt Møller reconstructs
%  Darwin's loss of orthodox Christian faith and argues that the
%  conceptual shift from a conventional God-concept to natural
%  selection was internally necessary.
% ============================================================
\documentclass[12pt,a4paper]{article}
\usepackage[utf8]{inputenc}
\usepackage[T1]{fontenc}
\usepackage[english]{babel}
\usepackage{microtype}
\usepackage{setspace}
\usepackage[margin=1.2in]{geometry}
\usepackage{fancyhdr}
\usepackage{hyperref}

\hypersetup{
  pdftitle={Darwin: His Life and His Doctrine (excerpt)},
  pdfauthor={J. P. Jacobsen and Vilhelm Møller},
  hidelinks
}

\pagestyle{fancy}
\fancyhf{}
\rhead{Møller, \emph{Darwin's Life} (1893), pp.\ 38--45}
\cfoot{\thepage}

\setstretch{1.3}

\title{\textbf{Darwin: His Life and His Doctrine}\\[0.4em]
\large Excerpt: ``Darwin's Loss of Faith''\\
(\emph{Darwin's Life}, pp.\ 38--45)}
\author{J.\ P.\ Jacobsen and Vilhelm Møller\\[4pt]
  \small Translated from the Danish}
\date{Gyldendal, Copenhagen, 1893}

\begin{document}
\maketitle
\thispagestyle{fancy}

\bigskip

\noindent Darwin does not believe that religious feeling was ever
strongly developed in him. Yet it was not infrequent and could
become vivid: he himself says that there were times when life
came into it. ``The state of mind,'' he says, ``that grand scenes
used to awaken in me, and which was very intimately connected
with a belief in God, differed in no essential respect from what
is usually called the feeling of the sublime''; and he compares
this state of mind with ``the powerful and somewhat indefinite
but analogous feelings excited by music.'' In other words, his
religious ideas were associated with, and sustained by, a
peculiar kind of poetry that was in him: his capacity for
enthusiasm, for being absorbed, seized, gliding into wondering,
admiring self-surrender.\footnote{See G.\ Brandes,
\emph{Naturalism in England}, p.\ 60~ff.}

``But now''---he adds, and this ``now'' would seem to refer
already to the final portion of the voyage---``now the most
sublime scenes could not produce in my mind any of those strong
and inwardly connected feelings with God.'' No: ``grand scenes''
simply no longer existed for him; he no longer experienced them
as grand. Only once more, in May of 1838, did he seem to
himself, during a church concert, to hear a choir that ``actually
made the walls shake''; otherwise, in all his letters from after
the voyage, one finds no preserved impression that is wholly
``sublime,'' only idyllic ones. A certain driving force in him
had, as noted, gone slack. Previously, under tremendous pressure,
it had driven his mood like a rainbow-glittering fountain jet
high, high into the blue of the ether: and on the topmost pearls
there had trembled in the sunlight a golden glimmer---a
star---the star that led wise men to the East. Now this driving
force, when set in motion, could only propel his mood a little
way into the air, and what had once seemed a heavenly star was
seen instead as a golden apple, wrought of human thought.

But what human beings have wrought, human beings reason about.
---``On board the \emph{Beagle} I was quite orthodox \ldots\ and
quoted the Bible as an unanswerable authority on points of
morality \ldots\ But at this time, from 1836 to 1839, I gradually
came to see that the Old Testament was no more to be trusted than
the sacred books of the Hindoos \ldots\ By further reflexion I
was gradually led to see that circumstances compelled me to
abandon my belief in Christianity as a divine revelation.''

In his mind he turned over possible proofs of the truth of the
Bible; but once Christianity no longer stood beyond refutation,
these artificial supports could not hold the conviction in
place. ``The rate of disbelief crept on very slowly,'' he
writes, ``and was at last complete. The disbelief came over me
without any stormlike suddenness or particular violence.''

With this there fell away as well a conception that, though
not logically necessary, was and is habitually bound up with
conventional Christian belief in God: the conception of God as
the direct producer of all the forms of organic life. In a
letter from somewhat later (1859), Darwin rightly observes
that just as the notion that thunder and lightning are the
effects of secondary causes has caused many to grieve---as
requiring them to give up the idea of direct divine
agency---so now many grieve to give up the idea of God as the
immediate author of the organic forms. For the phenomena of
the cosmos and the inorganic world, the majority of
nineteenth-century Christians had already ceased to believe that
God was the direct cause; but with respect to the phenomena of
organic life, they believed---and believe---that he is. The
conventional Christian belief in God was and is habitually bound
up with faith in the divine ``purposiveness'' of organic nature:
and with that no thoroughgoing hypothesis of evolution can be
reconciled. Only when Darwin had definitively departed---not
precisely from belief in a God, but from the conventional
God-concept---only then was it possible for him to conceive of
species as produced by the selection of chance variations
realized through the struggle for existence. For both with the
words ``struggle for existence'' and ``natural selection'' and
with the word ``chance,'' a long chain of natural causes is
inserted between the phenomena and God---with the word
``chance'' a chain so long that the human eye cannot trace more
than its uttermost link\ldots

Perhaps one further circumstance deserves to be drawn out in
explanation of the coming-to-be of ``Darwinism.'' There is no
direct statement in Darwin's letters and notebooks about the
emotional condition in which his rationalist critique was born.
But whenever he doubted the existence of God---and he never went
further than doubt, not even at his furthest extreme!---he
always supported himself upon this, and only this, objection:
``there is so much suffering in the world!'' That must have been
the most original and deepest motive for unbelief in his
consciousness.

His compassion for all that suffers likewise furnished him with
this weapon in the struggle against the doctrine of Providence.
His illness and his at times bitter revulsion at the injustice
of things sharpened the weapon. It must, of course, have been
most effective at those times when he ``looked darkly'' on life,
when his mood was such that the suffering of this world seemed
to him the prevailing thing, joy the vanishing exception.
But---had this mood fixed itself in him, had he, that is,
become a ``pessimist,'' he could not possibly have appropriated
and created the guiding idea that came to him through reading
Malthus's book. For it is not enough that the drive to
investigate is strong in a man, if he is to appropriate and
create guiding ideas from what he reads or encounters on his
way: there must also be, between what he sits with and what he
takes in, that mysterious affinity which chemistry, speaking of
substances, calls ``chemical kinship.'' And the idea he received
was a fortunate one---a deliberately optimistic idea.

He had mixed much with breeders and heard of their work: he
now conceived, to put it briefly, the struggle for existence as
the vast and patient ``breeder'' of the forms of organic life!
But already the very idea of ``breeding'' has something
festive about it. It calls to mind, it is true, toil and
disappointed hopes and failed attempts and sacrifices of many
kinds; but far, far more brightly does it illuminate the joy of
observation and experiment, and point towards the happily
attained goal: that a new variety is formed, useful and
satisfying to the one who formed it! In reality: behind the
Darwinian concept of ``natural selection'' (natural selection,
natural breeding) there glimmers a presentiment of Something or
Someone---``Nature,'' ``God,'' something Unfathomable and
Nameless!---that rejoices in the fact that the forms of life
change and are perfected to fit the shifting conditions of
existence. This concept carries with it a satisfaction in an
infinite and universal progress of utility\ldots\ And now with
respect to individuals---with respect to each single being in
the swarm of life that passes through the breeding-process of
the struggle for existence? Yes, for individuals as well---or
for the great preponderance of them---Darwin's idea of
``breeding'' presupposes a surplus of happiness. For the
breeding process would quickly come to an end if the bred
beings were more sick than healthy, felt more aversion than
appetite for food, were more unwilling than willing to
reproduce, and so forth.

One has, in fact, observed that just in the half-year during
which his central theory took shape, Darwin himself was, as it
were, in a state of deep and steady happiness. And it is easy to
guess the reason. It was in October 1838 that he read Malthus's
book,---and on the 29th of January 1839 his marriage took place
with his cousin, the admirable daughter of the admirable Josiah
Wedgwood, Emma Wedgwood! So far, then, has Eros had a hand in
the \emph{Origin of Species}. Without Darwin's happy love,
``Darwin's theory'' would not have come to be.

\end{document}
